\section{Results}
\label{sec:results}

We present results organized by retrieval quality (\S\ref{sec:main_quality}), per-repository analysis (\S\ref{sec:per_repo}), amortization comparison (\S\ref{sec:amortization_results}), cost analysis (\S\ref{sec:cost_results}), cross-benchmark validation (\S\ref{sec:cross_benchmark}), and a patching pilot (\S\ref{sec:patching_results}).

Results marked \textit{pending} indicate experiment cells that have not yet completed the readiness gates defined in \S\ref{sec:stats}. Completed cells report means over $\geq 3$ repeats with bootstrap 95\% CIs.

\subsection{Main Retrieval Quality}
\label{sec:main_quality}

Table~\ref{tab:main_results} presents the cross-repository macro F1 for all seven methods under the strict commit-fidelity track.

\begin{table*}[t]
\centering
\small
\begin{tabular}{lrrrrrrr}
\toprule
 & \multicolumn{4}{c}{SWE-bench repos} & \multicolumn{3}{c}{SWE-PolyBench repos} \\
\cmidrule(lr){2-5} \cmidrule(lr){6-8}
Method & Flask & Requests & Pytest & FastAPI & yt-dlp & LangChain & Keras \\
\midrule
\texttt{gm\_deterministic}$^{*}$ & 0.679 & 0.520 & 0.473 & 0.407$^\dagger$ & 0.297$^\ddagger$ & 0.452$^\ddagger$ & 0.274$^\ddagger$ \\
\texttt{gm\_progressive} & 0.803 & 0.700 & 0.550 & 0.417$^\dagger$ & 0.347$^\ddagger$ & 0.783$^\ddagger$ & 0.267$^\ddagger$ \\
\texttt{gm\_baseline} & 0.809 & 0.534 & 0.473 & 0.427$^\dagger$ & 0.198$^\ddagger$ & 0.312$^\ddagger$ & 0.253$^\ddagger$ \\
\texttt{rag\_progressive} & 0.704 & 0.733 & 0.602 & 0.421$^\dagger$ & 0.114$^\ddagger$ & 0.366$^\ddagger$ & 0.162$^\ddagger$ \\
\texttt{rag\_baseline} & 0.802 & 0.733 & 0.567 & 0.300$^\dagger$ & 0.300$^\ddagger$ & 0.433$^\ddagger$ & 0.140$^\ddagger$ \\
\texttt{raw\_rag\_function} & 0.340 & 0.256 & 0.242 & 0.255$^\dagger$ & 0.145$^\ddagger$ & 0.174$^\ddagger$ & 0.179$^\ddagger$ \\
\texttt{raw\_rag\_fixed} & 0.321 & 0.221 & 0.185 & 0.213$^\dagger$ & 0.149$^\ddagger$ & 0.132$^\ddagger$ & 0.152$^\ddagger$ \\
\bottomrule
\end{tabular}
\caption{Per-repository mean F1 under the strict commit-fidelity track. Flask, Requests, and Pytest: means over $\geq 3$ repeats for GM/RAG progressive/baseline; single run for deterministic and raw methods. $^\dagger$FastAPI: single run, SWE-bench ($n=10$). $^\ddagger$yt-dlp ($n=5$), LangChain ($n=10$), Keras ($n=10$): single run, SWE-PolyBench Verified. All single-run cells lack CIs; results are directional. $^{*}$Flask, Requests, and Pytest \texttt{gm\_deterministic} results use a tuned hyperparameter config (\texttt{configs/gm\_deterministic\_selected\_v1.json}, selected via Flask-only tuning with Requests/Pytest as holdout); FastAPI, yt-dlp, LangChain, and Keras use default weights. Scoring formula weights differ substantially between configurations (see Appendix~\ref{sec:scoring_details}).}
\label{tab:main_results}
\end{table*}

From the three completed SWE-bench repositories:
\begin{itemize}[leftmargin=1.5em]
  \item Graph methods win on Flask (\texttt{gm\_baseline}: 0.809), while RAG methods win on Requests and Pytest (\texttt{rag\_progressive}: 0.733 and 0.602 respectively).
  \item \texttt{gm\_deterministic} achieves competitive F1 (0.679, 0.520, 0.473) with zero LLM calls at query time --- all cost is setup embedding and lightweight query embedding.
  \item Non-LLM baselines (\texttt{raw\_rag\_*}) are not competitive, confirming that LLM-guided or graph-navigated retrieval provides substantial value.
  \item Cross-repo macro F1 across the three repos: RAG baseline 0.701, GM progressive 0.684, GM deterministic 0.557, RAG progressive 0.680, GM baseline 0.605.
\end{itemize}

\subsection{Per-Repository Paired Analysis (psf/requests)}
\label{sec:per_repo}

The most complete experiment cell is \texttt{psf/requests} with $n=10$ issues and repeated runs on both tracks. Table~\ref{tab:strict_requests} shows the strict-track repeat means.

\begin{table*}[t]
\centering
\begin{tabular}{lrrrr}
\toprule
Method & Mean F1 & Runtime Tokens & Setup Tokens & Total Cost \\
\midrule
GM (progressive) & 0.467 & 43,283 & 107,557 & 150,840 \\
GM (baseline) & 0.379 & 51,898 & 107,557 & 159,455 \\
RAG (progressive) & 0.286 & 15,587 & 414,082 & 429,669 \\
RAG (baseline) & 0.332 & 40,756 & 414,082 & 454,838 \\
\bottomrule
\end{tabular}
\caption{Strict commit-fidelity track: means over 3 repeats (\texttt{psf/requests}, $n=10$).}
\label{tab:strict_requests}
\end{table*}

% CLAIMS_LOCK: Keep quality phrasing as competitive/often-higher; cost-efficiency is the thesis. Do not frame as "significant advantage" without noting single-repo scope.
The paired F1 delta for GM-progressive minus RAG-progressive is $+0.181$ with bootstrap 95\% CI $[0.120, 0.233]$ on this repository. This indicates that GM-progressive is competitive with and often higher than RAG-progressive on \texttt{psf/requests}; cross-repo results are mixed (Table~\ref{tab:main_results}), consistent with the paper's cost-efficiency thesis rather than a quality-superiority claim.

\subsection{Strict vs.\ Same-Snapshot Amortization}
\label{sec:amortization_results}

Table~\ref{tab:amortization} compares the two evaluation tracks for the primary method pair.

\begin{table*}[t]
\centering
\begin{tabular}{llrrrr}
\toprule
Track & Method & Mean F1 & Runtime & Setup & Total Cost \\
\midrule
Strict & GM (progressive) & 0.467 & 43,283 & 107,557 & 150,840 \\
Strict & GM (deterministic) & 0.520 & 1,727 & 107,557 & 109,284 \\
Strict & RAG (progressive) & 0.286 & 15,587 & 414,082 & 429,669 \\
Same-snapshot & GM (progressive) & 0.544 & 42,136 & 9,018 & 51,154 \\
Same-snapshot & GM (deterministic) & 0.603 & 1,727 & 9,018 & 10,745 \\
Same-snapshot & RAG (progressive) & 0.259 & 15,345 & 36,736 & 52,081 \\
\bottomrule
\end{tabular}
\caption{Track-separated means (\texttt{psf/requests}, $n=10$; GM-progressive/RAG-progressive: means over $\geq 3$ repeats; GM-deterministic: single run). In same-snapshot mode, setup costs collapse due to single index build. GM-deterministic uses zero LLM runtime tokens.}
\label{tab:amortization}
\end{table*}

Key observations:
\begin{itemize}[leftmargin=1.5em]
  \item In the strict track, \texttt{commit\_repeat\_ratio = 0.0} and \texttt{cache\_hit\_rate = 0.0}: every issue requires a fresh index build.
  \item In the same-snapshot track, \texttt{commit\_repeat\_ratio = 0.9} and \texttt{cache\_hit\_rate = 0.9}: the graph is built once and reused.
  \item GM-progressive F1 improves from 0.467 (strict) to 0.544 (same-snapshot), suggesting that index stability benefits graph navigation.
  \item GM-deterministic same-snapshot cost is 10,745 tokens total vs.\ 109,284 strict---a 10$\times$ reduction from amortization---while F1 improves from 0.520 to 0.603. This illustrates the maximum cost benefit of world-model reuse with no LLM overhead.
  \item The same-snapshot GM-progressive minus RAG-progressive delta is $+0.286$ with 95\% CI $[0.219, 0.333]$.
\end{itemize}

Figure~\ref{fig:amortization} visualizes the strict vs.\ same-snapshot cost reduction alongside F1 for each method.

\begin{figure}[t]
\centering
\includegraphics[width=\linewidth]{fig_amortization}
\caption{Strict vs.\ same-snapshot total token cost (bars) and mean F1 (lines) for GM-progressive, GM-deterministic, and RAG-progressive on \texttt{psf/requests} ($n=10$). Cost reduction factors shown above bars.}
\label{fig:amortization}
\end{figure}

\subsection{Cost Analysis}
\label{sec:cost_results}

Table~\ref{tab:cost_cross_repo} presents the token cost decomposition across the three completed SWE-bench repositories (30 issues total, single run).

\begin{table*}[t]
\centering
\begin{tabular}{lrrr}
\toprule
Method & Runtime Tokens & Setup Tokens & Total Cost \\
\midrule
GM (deterministic) & 5,404 & 701,054 & 706,458 \\
GM (progressive) & 297,690 & 701,054 & 998,744 \\
GM (baseline) & 333,893 & 701,054 & 1,034,947 \\
RAG (progressive) & 123,317 & 2,268,044 & 2,391,361 \\
RAG (baseline) & 354,930 & 2,268,044 & 2,622,974 \\
Raw-func & 12,769 & 2,268,044 & 2,280,813 \\
Raw-fixed & 12,769 & 1,478,791 & 1,491,560 \\
\bottomrule
\end{tabular}
\caption{Token cost across 3 SWE-bench repos (Flask, Requests, Pytest; 30 issues total, single run per method). GM-deterministic has negligible runtime cost (no LLM calls at query time).}
\label{tab:cost_cross_repo}
\end{table*}

Figure~\ref{fig:f1_cost} plots the F1-vs-cost frontier across all seven methods for the three completed SWE-bench repositories.

\begin{figure*}[t]
\centering
\includegraphics[width=\linewidth]{fig_retrieval_f1_cost}
\caption{Retrieval F1 vs.\ per-issue cost (tokens, log scale) across seven methods for Flask, Requests, and Pytest. Graph methods (green) cluster at lower cost; RAG methods (red) at higher cost. Opacity encodes repository: Flask=100\%, Requests=80\%, Pytest=60\%.}
\label{fig:f1_cost}
\end{figure*}

GM methods achieve $\sim$3.2$\times$ lower setup cost than RAG methods (701K vs.\ 2.27M tokens) due to embedding compact graph metadata rather than full code bodies. GM-progressive total cost is 0.42$\times$ that of RAG-progressive despite 2.4$\times$ higher runtime token usage, because the setup cost dominates the total. GM-deterministic is the most cost-efficient retrieval method at 706K tokens total (0.30$\times$ of RAG-progressive), with zero LLM runtime spend.

\subsection{Cross-Benchmark Validation (SWE-PolyBench)}
\label{sec:cross_benchmark}

All single-run SWE-PolyBench experiments are now complete (Table~\ref{tab:main_results}, $^\ddagger$ columns). FastAPI results are from SWE-bench ($^\dagger$). All cross-benchmark cells are directional (single run, no CIs).

\textbf{FastAPI} (SWE-bench, $n=10$): GM-prog (0.417) and GM-base (0.427) are competitive with RAG-prog (0.421) at substantially lower setup cost (67K vs.\ 137K setup tokens).

\textbf{yt-dlp} (SWE-PolyBench, $n=5$): GM-prog (0.347) leads GM-det (0.297) and outperforms RAG-prog (0.114). RAG-baseline is competitive (0.300).

\textbf{LangChain} (SWE-PolyBench, $n=10$): GM-prog achieves the highest single-repo F1 across all evaluated repositories (0.783), substantially outperforming RAG-prog (0.366). This suggests graph-guided methods particularly benefit from structural connectivity in framework-scale codebases.

\textbf{Keras} (SWE-PolyBench, $n=10$): All methods obtain lower F1 than on SWE-bench repos, with GM-det (0.274) slightly ahead of GM-prog (0.267). RAG-baseline (0.140) trails raw-RAG methods (0.179/0.152), suggesting Keras issues involve less lexically predictable file locations. These lower values are consistent across methods and likely reflect benchmark difficulty rather than a method-specific effect.

\subsection{Patching Pilot (N=100, SWE-bench Verified)}
\label{sec:patching_results}

To evaluate end-to-end issue resolution, we run a patching pilot on 100 instances from the SWE-bench Verified subset across 9 repositories (see \S\ref{sec:patching_setup}). We report resolved rate (pass@1 under the SWE-bench Docker harness) and cost-per-resolved-issue. All results are labeled explicitly as exploratory: single run per method, no repeats, no CIs.

Table~\ref{tab:patching_resolved} shows per-repository resolved rates. Targeted reruns for sympy, sphinx-doc, and matplotlib (which had timeout confounding in the original runs) are complete; all numbers are final.

\begin{table*}[t]
\centering
\small
\begin{tabular}{lrrrr}
\toprule
Repository & Oracle & GM-prog & RAG-prog & None \\
\midrule
astropy/astropy ($n=12$) & 0.500 & 0.500 & 0.417 & 0.000 \\
matplotlib/matplotlib ($n=12$) & 0.167 & 0.250 & 0.083 & 0.083 \\
pallets/flask ($n=1$) & 1.000 & 1.000 & 0.000 & 0.000 \\
psf/requests ($n=8$) & 0.625 & 0.500 & 0.500 & 0.125 \\
pydata/xarray ($n=12$) & 0.500 & 0.417 & 0.583 & 0.000 \\
pytest-dev/pytest ($n=19$) & 0.368 & 0.526 & 0.421 & 0.053 \\
scikit-learn ($n=12$) & 0.500 & 0.500 & 0.417 & 0.000 \\
sphinx-doc/sphinx ($n=12$) & 0.333 & 0.167 & 0.250 & 0.000 \\
sympy/sympy ($n=12$) & 0.667 & 0.500 & 0.417 & 0.000 \\
\midrule
\textbf{Pooled} ($N=100$) & \textbf{0.450} & \textbf{0.430} & \textbf{0.380} & \textbf{0.030} \\
\bottomrule
\end{tabular}
\caption{Patching pilot resolved rates (pass@1) under SWE-bench Docker harness. All values final; targeted reruns for sympy, sphinx-doc, and matplotlib complete. Rows sorted alphabetically.}
\label{tab:patching_resolved}
\end{table*}

Table~\ref{tab:patching_cpr} presents the cost-per-resolved-issue (CPR) under two accounting views. The method-accounted view reflects design intent (GM builds graph only; RAG builds RAG index only). The as-run-consumed view reflects actual API spend during the current pilot runs, in which both graph and RAG indices were built for GM and RAG runs (dual-build confound, see \S\ref{sec:patching_setup}).

\begin{table*}[t]
\centering
\small
\begin{tabular}{lrrrrr}
\toprule
Method & Resolved & \multicolumn{2}{c}{CPR (tokens/resolved)} \\
\cmidrule(lr){3-4}
 & (out of 100) & Method-accounted & As-run-consumed \\
\midrule
Oracle & 45 & 33,594 & 33,594 \\
GM-progressive & 43 & 479,354 & 1,250,676 \\
RAG-progressive & 38 & 2,115,746 & 2,319,677 \\
None (cold-start) & 3 & 24,291 & 24,291 \\
\bottomrule
\end{tabular}
\caption{Cost-per-resolved-issue (CPR) for the patching pilot. Method-accounted CPR reflects design intent (GM builds graph only; RAG builds RAG index only); as-run-consumed reflects actual API spend including the dual-build confound (B10). None CPR appears low due to trivially low token spend on 3 cold-start resolutions. All numbers final.}
\label{tab:patching_cpr}
\end{table*}

Figure~\ref{fig:resolved_rates} visualizes the per-repository resolved rates across all four methods.

\begin{figure*}[t]
\centering
\includegraphics[width=\linewidth]{fig_resolved_rates}
\caption{Patching pilot resolved rates (pass@1) per repository and pooled, across oracle, GM-progressive, RAG-progressive, and cold-start (none) methods. Pooled row (N=100) shown in bold at bottom.}
\label{fig:resolved_rates}
\end{figure*}

\noindent Figure~\ref{fig:cpr} compares cost-per-resolved-issue under method-accounted and as-run-consumed accounting.

\begin{figure}[t]
\centering
\includegraphics[width=\linewidth]{fig_cpr_comparison}
\caption{Cost-per-resolved-issue (CPR) in tokens, log scale. Method-accounted (blue) reflects design intent; as-run-consumed (red) reflects pilot dual-build cost. The 4.4$\times$ ratio labels the GM vs.\ RAG method-accounted gap.}
\label{fig:cpr}
\end{figure}

Key observations (all exploratory, single-run pilot):
\begin{itemize}[leftmargin=1.5em]
  \item Oracle resolves 45\% of instances, providing the model ceiling under gold-file context. GM-progressive at 43\% is directionally close to oracle.
  \item GM-progressive resolves 5 more instances than RAG-progressive (43 vs.\ 38 out of 100). The quality difference is not statistically significant (McNemar exact $p=0.38$, $n_\text{discordant}=21$). Under method-accounted cost, GM-progressive CPR is approximately $4.4\times$ lower than RAG-progressive, consistent with the retrieval cost advantage.
  \item The GM-RAG gap narrows from 5pp all-in to 3.8pp when instances that exceeded the 480s wall-clock cap are excluded (timeout sensitivity analysis). A majority (16/21) of discordant instances are timeout-confounded; targeted reruns for sympy, sphinx-doc, and matplotlib are complete and reflected in the final numbers above.
  \item Cold-start (none) resolves only 3\% of instances, confirming that retrieval context provides substantial lift for the patch model.
\end{itemize}
